\documentclass[11pt]{article}
\usepackage[a4paper,margin=22mm]{geometry}
\usepackage[no-math]{fontspec}
\setmainfont{DejaVu Serif}
\usepackage{amsmath,amssymb,amsthm,mathtools,graphicx,xcolor,hyperref,xurl,xspace,ifthen}
\setlength{\emergencystretch}{4em}
\SetMathAlphabet{\mathbf}{normal}{OT1}{cmr}{bx}{n}
\SetMathAlphabet{\mathbf}{bold}{OT1}{cmr}{bx}{n}
\newtheorem{theo}{Théorème}[section]
\newtheorem{lemm}[theo]{Lemme}
\newtheorem{theorem}[theo]{Théorème}
\newtheorem*{remas*}{Remarques}
\newtheorem{remas}[theo]{Remarques}
\newtheorem*{exam*}{Exemple}
\newtheorem{ques}[theo]{Question}
\providecommand{\og}{«\nobreakspace}
\providecommand{\fg}{\nobreakspace»}
\newtheorem{lem}[theo]{Lemme}
\newtheorem{lemma}[theo]{Lemme}
\newtheorem{prop}[theo]{Proposition}
\newtheorem{coro}[theo]{Corollaire}
\newtheorem{corol}[theo]{Corollaire}
\newtheorem{conj}[theo]{Conjecture}
\newtheorem{Proposition}[theo]{Proposition}
\newtheorem{theoreme}[theo]{Théorème}
\theoremstyle{definition}
\newtheorem{defi}[theo]{Définition}
\newtheorem{exam}[theo]{Exemple}
\newtheorem{exem}[theo]{Exemple}
\newtheorem{rema}[theo]{Remarque}
\newtheorem{nota}[theo]{Notation}
\newtheorem*{rema*}{Remarque}
\newtheorem*{defi*}{Définition}
\newtheorem*{theo*}{Théorème}
\newtheorem*{lemm*}{Lemme}
\newtheorem*{prop*}{Proposition}
\newcommand{\enoncename}{}
\newtheorem{corpusenonce}[theo]{\enoncename}
\newenvironment{enonce}[2][]{\renewcommand{\enoncename}{#2}\begin{corpusenonce}}{\end{corpusenonce}}
\newenvironment{enonce*}[2][]{\par\medskip\noindent\textbf{#2.}\itshape}{\par\medskip}
\newenvironment{proofc}[1][\proofname]{\begin{proof}[#1]}{\end{proof}}
\newcommand{\Mc}{\mathcal{M}}
\newcommand{\pointir}{.\ ---\ }
\providecommand{\mathof}[1]{\mathrm{#1}}

\numberwithin{equation}{section}
\numberwithin{figure}{section}
\numberwithin{table}{section}
\def\endappendix{}
\tolerance=9999\relax
\renewcommand{\proofname}{Démonstration}
\renewcommand{\refname}{Bibliographie}
%~Mouliné par MaN_auto v.0.21.2 2019-03-29 18:08:38



\DeclarePairedDelimiterX{\intervalle}[2]{[}{]}{#1\mathclose{}\mathpunct{};#2}
\DeclarePairedDelimiterX{\intervalleoo}[2]{]}{[}{#1\mathclose{}\mathpunct{};#2}
\DeclarePairedDelimiterX{\intervalleof}[2]{]}{]}{#1\mathclose{}\mathpunct{};#2}
\DeclarePairedDelimiterX{\intervallefo}[2]{[}{[}{#1\mathclose{}\mathpunct{};#2}

\newcommand*\cA{{\mathcal A}}
\newcommand*\cF{{\mathcal F}}
\newcommand*\cS{{\mathcal S}}

\newcommand{\bbR}{{\mathbb{R}}}
\newcommand{\bbZ}{{\mathbb{Z}}}

\let\hat\widehat

\newenvironment{remasc*}{\let\pointir\relax \begin{remas*}}{\end{remas*}}
\newenvironment{theoc}{\let\pointir\relax \begin{theo}}{\end{theo}}
\newenvironment{step}{\begin{proof}}{\let\qed\relax \end{proof}}








%\DOI{10.NNNN/ahl.XXXX}









\graphicspath{{./figures/}}

\newcommand*{\mk}{\mkern -1mu}
\newcommand*{\Mk}{\mkern -2mu}
\newcommand*{\mK}{\mkern 1mu}
\newcommand*{\MK}{\mkern 2mu}

\hypersetup{urlcolor=purple, linkcolor=blue, citecolor=red}

\newcommand*{\romanenumi}{\renewcommand*{\theenumi}{\roman{enumi}}}
\newcommand*{\Romanenumi}{\renewcommand*{\theenumi}{\Roman{enumi}}}
\newcommand*{\alphenumi}{\renewcommand*{\theenumi}{\alph{enumi}}}
\newcommand*{\Alphenumi}{\renewcommand*{\theenumi}{\Alph{enumi}}}


\begin{document}
\section*{\raggedright 2683: Realization of every finite polynomial entropy greater than two by smooth Brouwer plane diffeomorphisms arising as time-one maps}
\noindent\textbf{Author:} Louis Hauseux; Frédéric Le Roux\par\noindent\textbf{Source:} Entropie polynomiale des homéomorphismes de Brouwer. Annales Henri Lebesgue 2 (2019), publisher version, DOI 10.5802/ahl.12. \url{https://ahl.centre-mersenne.org/item/AHL_2019__2__39_0/}
\par\noindent\textbf{License:} CC BY 4.0, \url{https://creativecommons.org/licenses/by/4.0/}. The original publisher PDF has the explicit license notice. See the saved source/license records.
\par\noindent\textbf{Editorial material note (not author text):} The selected problem is Theorem 1.1 part (3) for finite real alpha>2 only: a smooth Brouwer plane diffeomorphism arising as the time-one map of a smooth flow, whose one-point-compactification polynomial entropy is alpha. The complete original French author body retains all local orbit coding, entropy/localization, singular sets, finite-plane gluing, nested staircase plateaux and upper/lower word-count arguments. The alpha=infinity construction and other classifications are unselected. Original French language and theorem/proof headings are preserved; three original graph panels remain figure-only source assets with native caption and mathematical proof text. All proof text and formulas below are editable author TeX. Mathematical wording, language and proof order are retained. Only the article class, title matter and bibliography presentation are adapted for independent compilation; no proof has been generated or repaired.
\par\noindent\textbf{Editorial difficulty:} H1: The author realizes nonintegral entropy by carefully nested gluing plateaux and matches symbolic growth with singular-set localization. The interaction of planar topology, orbit coding and fractional complexity growth exceeds ordinary qualifying exercises.
\medskip\hrule\medskip
\noindent\textbf{Author statement, complete proof and local supporting text follow in source order.}
\section{Introduction}

L'entropie d'un système dynamique mesure le coefficient de croissance exponentielle du nombre d'orbites de longueur $n$ qu'on peut distinguer à une petite échelle. Quand cette entropie est nulle, c'est-à-dire lorsque la croissance du nombre d'orbites est sous-exponentielle, on peut essayer de mesurer un taux de croissance plus faible, ce qui conduit à définir différentes sortes d'entropies \og lentes \fg{} qui permettent de quantifier la complexité des systèmes d'entropie nulle. Cette idée a été réalisée sous diverses formes et dans des contextes variés. Pour des systèmes mesurables, Kushnirenko a introduit l'\emph{entropie séquentielle}, dans laquelle la croissance des orbites est examinée seulement pour une suite de temps adaptés, et l'a utilisée pour distinguer le flot horocyclique de son carré cartésien (\cite{Kus67}). Pour les actions mesurables des groupes moyennables, Anatole Katok et Jean-Paul Thouvenot ont défini l'\emph{entropie lente}, qui leur a permis de trouver des obstructions au problème de réalisation lisse (\cite{KatTho97}). Plus récemment, Adam Kanigowski a calculé la vitesse de croissance des orbites pour des flots avec un point fixe sur le tore, toujours d'un point de vue mesurable (\cite{Kan18}). En ce qui concerne les systèmes dynamiques topologiques, Jean-Pierre Marco a introduit l'entropie polynomiale dans le contexte des flots hamiltoniens intégrables (voir~\cite{Mar13}). Clémence Labrousse a étudié l'entropie polynomiale des homéomorphismes du cercle et des flots sur le tore. Elle a montré que les homéomorphismes du cercle ont une entropie polynomiale égale à $0$ ou $1$, la valeur $0$ caractérisant la classe de conjugaison la plus simple : celle des rotations (\cite{Lab13}). Elle s'est aussi intéressée à l'entropie polynomiale du flot géodésique sur le tore de dimension~2; C.~Labrousse et Patrick Bernard ont montré qu'un tel flot géodésique est d'entropie polynomiale égale à $1$ si et seulement si le tore est isométrique à un tore plat, dont la métrique est induite par celle du passage au quotient de $\mathbb R^2$ par un sous-groupe discret (\cite{BerLab16}). Dans un travail récent, Adam Kanigowski, Kurt Vinhage et Daren Wei calculent l'entropie polynomiale des flots quasi-unipotents sur un espace homogène de volume fini, et montrent dans ce contexte un principe du maximum qui relie l'entropie polynomiale topologique à l'entropie polynomiale métrique (\cite{KVW18}).

Nous nous sommes fixé deux objectifs pour cet article. Le premier objectif occupe toute notre section~\ref{s.wandering}; il ambitionne d'étudier, d'un point de vue très général, l'entropie polynomiale de la composante errante de n'importe quel système dynamique topologique; l'entropie polynomiale est ici particulièrement adaptée puisque nous montrerons que le taux de croissance du nombre d'orbites errantes est toujours polynomial (voir les remarques après la proposition~\ref{prop.localization}); l'entropie polynomiale peut alors être localisée autour d'un ensemble fini de points et calculée en faisant appel à la dynamique symbolique. Le deuxième objectif, qui occupe la section~\ref{s.brouwer}, consiste à appliquer nos résultats de la section précédente pour calculer l'entropie polynomiale de différents homéomorphismes de Brouwer; nous montrerons alors les résultats suivants :

\begin{theoc}\ \label{theo}
\begin{enumerate}%[topsep=0pt]
\item \label{1.1.1} L'entropie polynomiale d'un homéomorphisme de Brouwer $f$ est égale à $1$ si et seulement si $f$ est conjugué à une translation.

\item \label{1.1.2} Tout homéomorphisme de Brouwer qui n'est pas conjugué à une translation a une entropie polynomiale supérieure ou égale à $2$.

\item \label{1.1.3} Pour tout $\alpha \in \intervalle{2}{+\infty}$, on peut construire un homéomorphisme de Brouwer $f_{\alpha}$ d'entropie polynomiale $h_{pol}(f_{\alpha}) =\alpha$.
\end{enumerate}
\end{theoc}

Le dernier point est le plus intéressant ! En effet, il fournit des exemples naturels d'homéomorphismes ayant une entropie polynomiale non entière (en dynamique symbolique, une croissance polynomiale d'exposant non entier avait déjà été observée dans la complexité des systèmes Toeplitz par Julien Cassaigne et Juhani Karhumaki; voir~\cite[section~5]{CasKar97}). Il permet aussi de retrouver d'une nouvelle façon le résultat bien connu de l'indénombrabilité des classes de conjugaison des homéomorphismes de Brouwer, à l'aide d'un invariant numérique naturel. De plus, nos exemples ne sont pas particulièrement \og pathologiques\fg{} : ils sont tous le temps $1$ d'un flot $C^\infty$ constitué d'une seule composante de Reeb dont la frontière est un nombre fini de droites topologiques. Pour les construire, nous utiliserons une technique mise au point par Nakayama (\cite{Nak95}) puis reprise par François Béguin et l'un des auteurs (\cite{BegLer03}). L'entropie polynomiale des homéomorphismes de Brouwer semble d'ailleurs être assez fortement reliée à la notion d'\emph{ensembles oscillants} (autre invariant de conjugaison dont il est question dans~\cite{BegLer03}); en particulier, il est fort possible qu'une entropie polynomiale élevée engendre nécessairement un ensemble oscillant non vide.


\section{Entropie polynomiale de systèmes dynamiques errants}\label{s.wandering}

\subsection{Contexte et définition}

Par soucis de simplicité, nous considérons un espace métrique $W$ et un homéomorphisme $g: W \to W$, quoique tout ce qui suit puisse probablement être généralisé aux applications continues non inversibles. Rappelons qu'un ensemble $Y$ est \emph{errant} si $g^n(Y) \cap Y =\emptyset$ pour tout $n \neq 0$ et qu'un point est errant s'il admet un voisinage errant. On désignera l'ensemble des points non-errants de $g$ par $NW(g)$ (pour \emph{Non-Wandering}). Cet ensemble est invariant par $g$. Soit $\hat W$ le résultat du passage au quotient de $W$ par $NW(g)$, pour lequel tous les éléments de $NW(g)$ sont identifiés. Soit alors $\hat g: \hat W \to \hat W$ l'homéomorphisme induit par $g$. Remarquons que tout point de $\hat W$ est errant pour $\hat g$ excepté un : celui qui est l'image de $NW(g)$; on notera ce point $\infty$. Le théorème de récurrence de Poincaré affirme que la seule mesure qui est invariante pour $\hat g$ est la mesure de Dirac au point $\infty$, et en vertu du principe variationnel, l'entropie topologique de $\hat g$ est nulle. On est ainsi naturellement amené à la définition suivante.

\begin{defi}
L'\emph{entropie polynomiale errante} de $g$ est l'entropie polynomiale de $\hat g: \hat W \to \hat W$.
\end{defi}

L'espace topologique $\hat W$ peut être muni d'une distance : une distance naturelle $\delta$ est donnée, par exemple, par les formules
\begin{gather*}
\delta(x,NW(g))= \inf\{d(x,y), y \in NW(g)\} \quad\text{si } x \not \in NW(g)
\\
\text{et}\quad \delta(x,y) = \min(d(x,y), \delta(x,NW(g))+\delta(y,NW(g)))) \quad\text{si } x,y \not \in NW(g).
\end{gather*}
On se limite donc à étudier l'entropie polynomiale des homéomorphismes d'un espace métrique compact dont l'ensemble non-errant est réduit à un unique point : c'est le cadre que nous adoptons pour notre étude.

%\bigskip

Tout au long des pages qui vont suivre, $X$ désignera un espace métrique compact et $\infty$ un point de $X$. On considère un homéomorphisme $f:X \to X$ dont $\infty$ est un point fixe. De plus, $f$ vérifiera l'hypothèse suivante : \emph{tout point de $X$ excepté $\infty$ est errant}.

Rappelons ici la définition de l'entropie polynomiale de $f$ (voir~\cite{Mar13}). \'Etant donnés un réel $\varepsilon>0$ et un entier $n > 0$, un ensemble $E$ de points de $X$ est dit $(n,\varepsilon)$-séparé si, pour toute paire $x,y \in E$, il existe un entier $k \in \{0, \dots, n-1\}$ tel que : $d(f^k(x), f^k(y))>\varepsilon$. Par compacité de $X$, $E$ est nécessairement fini; on note $S(n,\varepsilon)$ le nombre maximal d'éléments de tels ensembles $(n,\varepsilon)$-séparés. L'entropie polynomiale de $f$ est alors définie par :
\[
h_{pol}(f) = \lim_{\varepsilon \to 0} \limsup_{n \to +\infty} \frac{\log S(n,\varepsilon)}{\log (n)}.
\]
D'une façon alternative, on aurait pu tout aussi bien définir cette quantité
à l'aide des ensembles $(n,\varepsilon)$-couvrants ou encore à l'aide des recouvrements dynamiques (les arguments sont identiques à ceux déployés pour l'entropie classique, voir par exemple le chapitre dédié à l'entropie topologique dans~\cite{Wal82}). En particulier, c'est un invariant de conjugaison topologique (cf.~\cite{Mar13}).

\subsection{Encodage}\label{ss.coding}

Soit $\cF$ une famille finie de parties non-vides de $X \setminus \{\infty\}$. On désigne par $\cup \cF$ l'union de tous les éléments de $\cF$, et par $\infty_\cF$ le complémentaire de $\cup \cF$, qu'on notera simplement $\infty$ quand il n'y aura pas de risques de confusion avec le point $\infty$. Fixons un entier $n>0$. Soient $\underline{x} = (x_{0}, \dots, x_{n-1})$ une suite de $n$ points de $X$ et $\underline{w} = (w_0, \dots, w_{n-1})$ une suite de $n$ éléments de $\cF\cup\{\infty_\cF\}$. On dira alors que $\underline{w}$ est un \emph{encodage} de $\underline{x}$ par rapport à $\cF$ si $x_{k} \in w_k$ pour tout entier $k \in \{0, \dots, n-1\}$. Remarquons que si les éléments de $\cF$ ne sont pas disjoints, un même élément $\underline{x}$ peut avoir plusieurs encodages. On désigne par $A_{n}(\cF)$ l'ensemble de tous les encodages possibles des orbites $(x, f(x), \dots, f^{n-1}(x))$ de longueur $n$. On peut alors définir \emph{l'entropie polynomiale de la famille $\cF$} comme étant :
\[
h_{pol}(f \,; \cF) = \limsup_{n \to \infty} \frac{\log \sharp A_{n}(\cF)}{\log (n)}.
\]
Si $\cF = \{Y\}$ est réduit à un seul élément, alors on écrira plus simplement $h_{pol}(f; Y)$ à la place de $h_{pol}(f; \{Y\})$. On omettra souvent aussi de préciser l'homéomorphisme \og $f$\fg{} quand le contexte sera sans ambigu\"{\i}tés\footnote{Attention cependant : cette notation a déjà été utilisée pour désigner autre chose (cf.~\cite[définition~1]{Mar13}).}.

\begin{exam*}
Soit $A$ {l'application linéaire} du plan $\mathbb R^2$ : $(x,y) \mapsto (2x, y/2)$. Pour répondre aux hypothèses précédemment requises, on commence par compactifier le plan en lui ajoutant un point à l'infini qu'on identifie au point $0$; on obtient alors notre espace $X$ et notre homéomorphisme $f$, qui admet $\infty$ comme seul point fixe. Comme l'on travaille ensuite avec des sous-ensembles compacts de $X\setminus \{\infty\}$, on pourra les assimiler à des sous-ensembles de $\mathbb R ^2 \setminus \{0\}$. Soient $Y_1, Y_2$ deux disques ne contenant pas l'origine et dont l'intérieur intersecte respectivement l'axe des ordonnées et l'axe des abscisses. Par souci de simplicité, on supposera ces disques assez petits de sorte qu'ils soient disjoints de leur image. Tout d'abord, on a : $h_{pol}(A;Y_1)=h_{pol}(A;Y_2)=1$. En effet, les éléments, par exemple, de $\cA_n(Y_1)$ sont précisément tous les mots de la forme $(\infty, \dots, \infty, Y_1, \infty, \dots, \infty)$ et sont donc au nombre de $n$. Ensuite, on a $h_{pol}(A;\{Y_1, Y_2\})=2$. En effet, un élément de $\cA_n(\{Y_1, Y_2\})$ est un mot de la forme $(\infty, \dots, \infty, Y_1, \infty, \dots, \infty, Y_2, \infty, \dots, \infty)$; il y en a donc au plus $n(n-1)$\dots ce qui fournit la borne supérieure. Pour la borne inférieure, on s'aperçoit qu'il existe un entier $L$ tel que pour tout $n \geq L$, $f^n(Y_1)$ intersecte $Y_2$; ainsi, $\cA_n(\{Y_1, Y_2\})$ contient tous les mots de la forme précédente où $Y_1$ et $Y_2$ sont séparés d'une distance $k \geq L-1$. On obtient alors l'équivalent voulu :
\[
\sharp \cA_n(\{Y_1, Y_2\}) \geq \sum_{k=L-1}^{n-2} (n-k-1) \sim \sum_{\ell=1}^n \ell \sim \frac{1}{2} n^2.
\]
On verra une généralisation de cet exemple grâce à la notion d'\textit{ensemble singulier} (se reporter à la section~\ref{ss.singular-set} plus bas).
\end{exam*}

%\bigskip


Pour tout sous-ensemble $Y$ de $X\setminus \{\infty\}$, on note $M(Y)$ le nombre maximal de fois qu'une orbite peut passer dans $Y$:
\[
M(Y) = \sup_{x \in X} \sharp \{n, f^n(x) \in Y \}.
\]
Comme $Y$ est compact, on peut le recouvrir par un nombre fini d'ouverts errants; or chaque orbite passe au plus une fois dans chacun de ces ouverts, on en déduit que $M(Y) < +\infty$. Considérons une famille finie $\cF$ de parties de $X\setminus \{\infty\}$ telle que $M(\cup \cF) < +\infty$. Si $n$ est un entier beaucoup plus grand que $M(\cup \cF)$, alors tout mot de $\cA_n(\cF)$ est constitué pour une grande majorité de la lettre $\infty_\cF$. Cette remarque conduit au résultat suivant.

\begin{lemm}\label{lemm.add-mon}
Considérons une famille finie $\cF$ de sous-ensembles de $X$ tels que $M(\cup\cF) < +\infty$. On a les propriétés suivantes :
\begin{enumerate}
\item \label{2.2.1} (décroissance) si $\cF'$ est une autre famille finie de sous-ensembles de $X$ telle que chacun de ses éléments soit inclus dans un élément de $\cF$, alors :
\[
h_{pol}(\cF') \leq h_{pol}(\cF),
\]
\item \label{2.2.2} (union)
\[
h_{pol}(\cup \cF) = h_{pol}(\cF).
\]
\item \label{2.2.3} (restriction) si $\cF=\{Y_1, \dots, Y_L \}$ possède deux éléments $Y_1$ et $Y_2$ dont l'union $Y_{1} \cup Y_2$ est errante, alors :
\[
h_{pol}(\cF) = \max(h_{pol}(Y_{1},Y_3, \dots, Y_{L}), h_{pol}(Y_{2},Y_3, \dots, Y_{L})).
\]
\end{enumerate}
\end{lemm}

\begin{proof}
Pour montrer le premier point (la {décroissance}), posons un entier $n > 0$ et regardons une application $\Phi$ de $\cA_n(\cF')$ dans $\cA_n(\cF)$ définie ainsi : étant donné un mot $w'$ de $\cA_n(\cF')$, on choisit un représentant $x$ tel que $w'$ soit l'encodage de $(x, \dots, f^{n-1}(x))$ relatif à $\cF'$, on donne à $\Phi(w')$ la valeur d'un encodage de $(x, \dots, f^{n-1}(x))$ pour $\cF$. \`A présent réfléchissons : combien un mot {donné} $w \in \cA_n(\cF)$ peut-il avoir d'antécédents {$w' \in \cA_n(\cF')$} ? ou, dit autrement, combien y a-t-il de possibilités d'encodages $w'$ relatifs à $\cF'$ pour l'orbite de $x$ ? Si la $k$-ième lettre de $w$ est \og$\infty$\fg, cela signifie que $f^k(x) \not \in \cup \cF$, auquel cas $f^k(x) \not \in \cup \cF'$ et la $k$-ième lettre de $w'$ est aussi nécessairement \og$\infty$\fg. De plus, $w$ diffère du mot $(\infty, \dots, \infty)$ pour au plus $M(\cup \cF)$ lettres. Puisque $w'$ est un mot sur un alphabet comportant $\sharp \cF'+1$ lettres, cela ne laisse plus que
\[
C = \left(\sharp \cF'+1 \right)^{M(\cup \cF)}
\]
possibilités pour $w'$. On en déduit que :
\[
\sharp \cA_n(\cF') \leq C \times \sharp \cA_n(\cF).
\]
Comme $C$ ne dépend pas de $n$, on obtient l'inégalité voulue : $h_{pol}(\cF') \leq h_{pol}(\cF)$.

%\bigskip

Penchons-nous à présent sur le second point. Grâce à la relation précédente appliquée aux familles $\cF$ et $\{\cup \cF\}$, on possède déjà une première inégalité :
\[
h_{pol}(\cF) \leq h_{pol}(\cup \cF).
\]
Le sens réciproque vient de l'inégalité très simple suivante (car $\cA_n(\cF)$ permet de faire les mêmes encodages que $\cA_n(\cup \cF)$ avec plus de lettres à disposition) :
\[
\sharp \cA_n(\cup \cF) \leq \sharp \cA_n(\cF).
\]

%\bigskip

Enfin, examinons le troisième point. En appliquant le premier point deux fois, on a déjà que :
\[
h_{pol}(\cF) \geq \max(h_{pol}(Y_{1},Y_3, \dots, Y_{L}), h_{pol}(Y_{2},Y_3, \dots, Y_{L})).
\]
Prouvons maintenant l'inégalité réciproque. Comme $Y_1\cup Y_2$ est errant, aucun mot de $\cA_n(\cF)$ ne peut contenir à la fois les lettres \og$Y_1$\fg{} et \og$Y_2$\fg; en conséquence de quoi
\[
\cA_n(\cF) \subseteq \cA_n(Y_1, Y_3, \dots, Y_L) \cup \cA_n(Y_2, Y_3, \dots, Y_L).
\]
D'où
\begin{align*}
\sharp A_{n}(\cF) & \leq \sharp \cA_n(Y_1, Y_3, \dots, Y_L) +\sharp \cA_n(Y_2, Y_3, \dots, Y_L) \\
& \leq 2 \max \left(\sharp \cA_n(Y_1, Y_3, \dots, Y_L), \sharp \cA_n(Y_2, Y_3, \dots, Y_L))\right)
\end{align*}
et on obtient l'inégalité souhaitée.
\end{proof}

\subsection{Localisation}\label{ss.localization}

Soient $x_{1}, \dots, x_{L}$ des points de $X \setminus \{\infty\}$. Choisissons pour chaque $\ell \in \{1,\dots,L\}$ une suite d'ensembles décroissante $(U_{\ell, n})_{n \geq 0}$ qui forme une base de voisinage de $x_{\ell}$. Par décroissance (premier point de notre précédent lemme), la suite
\[
\left(h_{pol}(U_{1,n},U_{2,n}, \dots, U_{L,n})\right)_{n \geq 0}
\]
est décroissante. On appelle sa limite \emph{entropie polynomiale locale aux points $(x_{1}, \dots, \allowbreak x_{L})$} et on la note $h^{loc}_{pol}(f;x_{1}, \dots, x_{L})$ (on omettra souvent \og$f$\fg). De plus, {toujours grâce à la propriété de décroissance,} cette limite ne dépend que du choix des $x_{\ell}$ et non pas de leur base de voisinages. Si l'on nomme $\cS = \{x_{1}, \dots, x_{L}\}$ la famille de points, on pourra noter $h_{pol}(x_{1}, \dots, x_{L})$ simplement $h_{pol}(\cS)$.


\begin{prop}[Localisation]\label{prop.localization}
L'entropie polynomiale de $f$ est donnée par les formules suivantes :
\begin{align*}
h_{pol}(f) & = \sup\left\{h^{loc}_{pol}(f;\cS)\right\} \\
&= \sup \left\{ h_{pol}(f;Y)\right\}
\end{align*}
où le premier supremum porte sur toutes les familles finies $\cS$ de points de $X\setminus \{\infty\}$ et le second supremum sur {toutes les parties compactes} $Y$ de $X\setminus \{\infty\}$.
\end{prop}

\begin{remasc*}\ 
\begin{enumerate}
\item Si $\cF=\{Y\}$ et $Y$ est errant, alors, pour les mêmes raisons que dans l'exemple de la section~\ref{ss.coding}, on a que : $\sharp \cA_n(\cF)= n$ et par suite $h_{pol}(f; \cF)=1$. On en déduit, grâce à la proposition, que l'entropie polynomiale de $f$ est toujours plus grande que $1$.
\item On a aussi la majoration suivante : $h_{pol}(Y)\leq M(Y)$. En effet, l'encodage d'un mot de $\cA_n(Y)$ fait apparaître au plus $M(Y)$ fois la lettre \og$Y$\fg, le cardinal de $\cA_n(Y)$ est donc borné par $n^{M(Y)}$. Au cours de la preuve de la Proposition~\ref{prop.localization}, on verra que pour tout $\varepsilon>0$, la croissance de $S(n,\varepsilon)$ est encore au plus polynomiale. On en conclut qu'en dynamique errante, la croissance du nombre d'orbites est au moins linéaire et au plus polynomiale.
\item Si $\cS$ est un ensemble fini, on peut trouver une famille $\{U_x, x \in \cS\}$ de voisinages errants des points de $\cS$. On a alors :
\[
h_{pol}(\{U_x, x \in \cS\})=h_{pol}\left(\bigcup_{x \in \cS} U_x\right) \leq M\left(\bigcup_{x \in \cS} U_x\right) \leq \sharp \cS
\]
et l'on en déduit que l'entropie polynomiale locale {aux points} de $\cS$ est {au plus égale} au nombre d'éléments de $\cS$.
\end{enumerate}
\end{remasc*}

%\bigskip

Le lemme suivant fournit un minorant de l'entropie polynomiale (on aura besoin d'un sous-lemme pour le montrer).

\begin{lemm}\label{lem.inequality}
Soit $\cF$ une famille finie de sous-ensembles compacts de $X \setminus\{\infty\}$. Alors
\[
h_{pol}(f;\cF) \leq h_{pol}(f).
\]
\end{lemm}

\begin{enonce}{Sous-Lemme}\label{sublemma.cut1}
Pour tout sous-ensemble compact $Y$ de $X \setminus \{\infty\}$ et tout $\varepsilon>0$, il existe un entier $L$ et $\{Y_1, \dots,Y_L\}$ une famille de $L$ sous-ensembles compacts de $Y$, {errants}, disjoints deux à deux et de diamètres inférieurs à $\varepsilon$ tels que :
\[
h_{pol}(f;\{ Y_1, \dots,Y_L \}) = h_{pol}(f;Y).
\]
\end{enonce}

\begin{proof}[Preuve du sous-lemme]
Par hypothèse, on sait que tout point de $Y$ admet un voisinage errant. Par compacité de $Y$, on peut trouver un sous-recouvrement errant de $Y$. Le lemme de Lebesgue nous permet alors, quitte à diminuer $\varepsilon$, de supposer que tout sous-ensemble de $Y$ de diamètre plus petit que $2\varepsilon$ soit errant. En recourant encore une fois à la compacité de $Y$, on peut donc recouvrir $Y$ par une famille $\{ Y_1, \dots,Y_L\}$ {de parties compactes de $Y$} de diamètres inférieurs à $\varepsilon$. {En particulier,} pour tous $i,j \in \{1, \dots, L\}$, si $Y_i$ intersecte $Y_j$ alors $Y_i \cup Y_j$ est encore un ensemble errant.

On va appliquer maintenant l'algorithme suivant pour construire une suite décroissante $\cF_0, \dots, \cF_n$ de familles {de parties} de $Y$ qui ont toutes la même entropie polynomiale. Initialisons l'algorithme avec $\cF_0 =
\{ Y_1, \dots,Y_L\}$. \`A l'étape $k$, on a construit une famille $\cF_k$. Si ses éléments sont deux à deux disjoints, on s'arrête. Sinon, on en choisit deux $Y_i, Y_j \in \cF_k$ qui s'intersectent. Par construction, leur union $Y_i \cup Y_j$ est errante et, en appliquant le dernier point du lemme~\ref{lemm.add-mon} (celui sur la restriction d'une famille quand l'union de deux de ses éléments est errante), on choisit selon le cas $\cF_{k+1} = \cF_k \setminus \{ Y_i\}$ ou $\cF_{k+1} = \cF_k \setminus \{ Y_j\}$ pour avoir
\[
h_{pol}(f;\cF_{k+1}) = h_{pol}(f;\cF_{k}).
\]

L'algorithme retourne une sous-famille $\cF_n = \{Y_{i_1}, \dots, Y_{i_{L-n}}\}$ de sous-ensembles compacts de $Y$ qui sont de diamètres inférieurs à $\varepsilon$ et deux à deux disjoints; elle vérifie de plus :
\[
h_{pol}(f;\cF_n) = h_{pol}(f;Y).\qedhere
\]
\end{proof}

\begin{proof}[Preuve du lemme]
Par la propriété d'union du lemme~\ref{lemm.add-mon}, on a : $h_{pol}(\cup \cF) = h_{pol}(\cF)$; il nous suffit donc de montrer que $h_{pol}(Y) \leq h_{pol}(f)$ quand $Y$ est un sous-ensemble compact de $X \setminus\{\infty\}$. D'après le sous-lemme qu'on vient de montrer, il existe une famille $\cF$ composée d'ensembles errants deux à deux disjoints et dont l'entropie polynomiale est égale à celle de $Y$. Il ne reste donc plus qu'à montrer que
\[
h_{pol}(f;\cF) \leq h_{pol}(f).
\]
Soient $U_1, \dots, U_L$ des voisinages errants respectivement de $Y_1, \dots, Y_L \in \cF$, deux à deux disjoints. Soit $\varepsilon>0$ plus petit que chacune (pour $\ell=1, \dots, L$) des distances des $Y_\ell$ au complémentaire de $U_\ell$. Fixons un entier $n > 0$. Pour toute sous-famille $\cF' \subset \cF$, soit $\cA_n(\cF,\cF')$ le sous-ensemble de $\cA_n(\cF)$ contenant tous les éléments dont l'encodage utilise exactement toutes les lettres de $\cF' \cup \{\infty_\cF\}$. Si l'on considère maintenant deux points $x,y$ de $X$ et deux mots $\underline{w} = (w_{0}, \dots, w_{n-1})$ et $\underline{z} = (z_{0}, \dots, z_{n-1})$ de $\cA_n(\cF,\cF')$ qui encodent respectivement les orbites $(x, \dots, f^{n-1}(x))$ et $(y, \dots, f^{n-1}(y))$, alors si les mots $\underline{w}$ et $\underline{z}$ sont distincts, les points $x$ et $y$ sont $(n,\varepsilon)$-séparés. En effet, soit $i \in \{0, \dots, n-1\}$ tel que $w_i \neq z_i$. Si l'on a en même temps $w_i \neq \infty, z_i \neq \infty$, c'est que $f^i(x)$ et $f^i(y)$ appartiennent à deux ensembles $Y_{\alpha} \neq Y_\beta$ et sont donc séparés d'une distance plus grande que $\varepsilon$. Si, disons, $w_i=\infty$ alors $f^i(y) \in Y_{z_i}$ et $f^i(x) \not \in Y_{z_i}$. Or, par définition de $\cA_n(\cF,\cF')$, il existe $j \neq i$ dans $\{0, \dots, n-1\}$ tel que $f^j(x) \in Y_{z_i} \subset U_{z_i}$. Comme $U_{z_i}$ est un ensemble errant, on obtient que $f^i(x) \not \in U_{z_i}$, encore une fois $f^i(x)$ et $f^i(y)$ sont à plus de $\varepsilon$ {: finalement} $x$ et $y$ sont bien $(n,\varepsilon)$-séparés. En conséquence, $\sharp \cA_n(\cF,\cF') \leq S(n, \varepsilon)$. De plus, les $2^L$ ensembles $\cA_n(\cF,\cF')$ forment une partition de $\cA_n(\cF)$; ainsi :
\[
\sharp \cA_n(\cF) \leq 2^L S(n, \varepsilon)
\]
et puis
\[
h_{pol}(f; \cF) \leq \limsup_{n \to \infty} \frac{\log S(n, \varepsilon)}{\log (n)}.
\]
Cette inégalité étant valide pour tout $\varepsilon$ suffisamment petit, on obtient finalement que
\[
h_{pol}(f; \cF) \leq h_{pol}(f).\qedhere
\]
\end{proof}

\begin{proof}[Preuve de la proposition~\ref{prop.localization}]
Le lemme ci-dessus prouve déjà que $h_{pol}(f)$ est plus grand que les deux supremums apparaissant dans la proposition. Voyons à présent les inégalité réciproques et d'abord celle portant sur les ensembles compacts ne contenant pas $\infty$. Soit $\varepsilon >0$. Soit $\{Y_1, \dots, Y_L\}$ une famille de sous-ensembles (à priori non disjoints) de $X \setminus \{\infty\}$ de diamètres inférieurs à $\varepsilon$ et tels que $Y_\infty = X \setminus (Y_1 \cup \dots \cup Y_L)$ soit aussi de diamètre plus petit que $\varepsilon$. Soit un entier $n >0$. Considérons à présent un ensemble $E$ qui soit $(n,\varepsilon$)-séparé. Pour chaque point $x$ de $E$, on choisit un encodage $\underline{\alpha}(x)$ de l'orbite $(x, f(x), \dots, f^{n-1}(x))$ sur l'alphabet $Y_\infty \cup \{Y_1, \dots, Y_L\}$. Comme $E$ est $(n,\varepsilon$)-séparé et que les ensembles $Y_1, \dots, Y_L, Y_\infty$ sont de diamètres plus petits que $\varepsilon$, une application d'encodage $\underline{\alpha}$ de $E$ dans l'ensemble des mots est forcément injective. Ainsi
\[
S(n,\varepsilon) \leq \sharp \cA_n(\{Y_1, \dots, Y_L\}).
\]
En divisant par $\log(n)$ et en faisant tendre $n$ vers l'infini, on obtient :
\[
\limsup_{n \to \infty} \frac{\log S(n,\varepsilon)}{\log (n)} \leq h_{pol}(Y_{1}, \dots, Y_{L}) = h_{pol}(Y)
\]
en posant $Y=Y_{1} \cup \dots \cup Y_{L}$ (d'après la propriété d'union). Maintenant, c'est $\varepsilon$ qu'on fait tendre vers $0$ pour avoir :
\[
h_{pol}(f) \leq \sup \{ h_{pol}(Y), Y \text{ partie compacte de } X\setminus \{\infty\} \}
\]
et prouver ainsi la première égalité.
%\bigskip

Considérons à présent la seconde égalité. Soit $\alpha>0$. On veut trouver $L$ points $x_1, \dots x_L$ dont l'entropie polynomiale soit plus grande que $h_{pol}(f)-\alpha$. D'après ce que l'on vient de voir, on peut déjà trouver un ensemble compact $Y$ ne contenant pas $\infty$ dont l'entropie polynomiale soit plus grande que $h_{pol}(f)-\alpha$. En utilisant le sous-lemme~\ref{sublemma.cut1}, on sait l'existence d'une famille $Y_1, \dots, Y_L$ de parties compactes errantes deux à deux disjointes dont l'entropie polynomiale égale celle de $Y$.

\begin{enonce}{Sous-Lemme}\label{sublemma.cut2}
Pour tout $\eta>0$, il existe $Y'_1, \dots, Y'_L$ inclus respectivement dans $Y_1, \dots, Y_L$, dont les diamètres sont plus petits que $\eta$ et vérifiant :
\[
h_{pol}(f; Y'_{1}, \dots, Y'_{L})= h_{pol}(f; Y_{1}, \dots, Y_{L}).
\]
\end{enonce}

\begin{proof}
On découpe les $Y_i$ (comme précédemment) en de plus petits morceaux et l'on recourt à la propriété de restriction.
\end{proof}

Ce sous-lemme permet de construire une suite décroissante $(Y_i^n)_{n \geq 0}$, $i=1 \dots L$ avec $Y_i^0 = Y_i$ et dont les éléments ont un diamètre qui tend vers $0$ et qui laissent inchangés $h_{pol}(f; Y_{1}^n, \dots, Y_{L}^n)$ quel que soit $n$; {ce nombre reste par conséquent plus grand que $h_{pol}(f)-\alpha$}. {Notons} $x_1, \dots x_L$ les limites respectives de $Y_1, \dots, Y_n$ :
\[
\{x_i\} = \bigcap_n Y_i^n,
\]
alors, pour tous voisinages $U_1, \dots, U_L$ de $x_1, \dots, x_L$, on peut trouver $n$ assez grand pour que $Y_i^n \subset U_i$ pour n'importe quel $i$. La propriété de décroissance nous donne :
\[
h_{pol}(f; U_{1}, \dots, U_{L}) \geq h_{pol}(f; Y_{1}^n, \dots, Y_{L}^n) > h_{pol}(f)-\alpha.
\]
Finalement, $h^{loc}_{pol}(f; x_{1}, \dots, x_{L})$ est supérieur à $h_{pol}(f)-\alpha$, {ce qui termine la\linebreak preuve}.
\end{proof}

Terminons cette section avec le résultat pratique suivant sur la localisation de l'entropie.

\begin{lemm}\label{lemm.iterates}
Pour tous points $x_1, \dots, x_L$ de $X \setminus\{\infty\}$ et tous entiers $n_{1}, \dots, n_{L}$,
\[
h^{loc}_{pol}(f;x_{1}, \dots, x_{L}) = h^{loc}_{pol}(f;f^{n_{1}}x_{1}, \dots, f^{n_{L}}x_{L}).
\]
\end{lemm}

\begin{proof}
Commençons par montrer la formule préliminaire suivante~:
\[
h^{loc}_{pol}(f;x_1, f^{n_1}(x_{1}), x_2, \dots, x_{L}) = h^{loc}_{pol}(f;x_{1}, \dots, x_{L}).
\]
Par décroissance, on sait déjà que le membre de droite est plus petit que celui de gauche. Si $f^{n_1}(x_1)$ est égal à l'un des $x_j$, l'égalité est claire. De même, on peut supposer sans perte de généralité que les $x_j$ sont tous distincts. Prenons des voisinages $U_1, f^{n_1}(U_1), U_2, \dots U_L$ deux à deux disjoints de $x_1, f^{n_1}(x_{1}), x_2, \dots, x_{L}$. Construisons une application
\[
\Phi: \cA_{n}(U_1, f^{n_1}(U_1), \dots, U_L) \to \cA_{n+n_1}(U_1, \dots, U_L)
\]
de la façon suivante. \'Etant donné un mot $w \in \cA_{n}(U_1, f^{n_1}(U_1), \dots, U_L)$, on peut choisir $x$ tel que $w$ soit l'encodage de $x, \dots, f^{n-1}(x)$. On définit alors $\Phi(w)$ comme étant l'encodage de $f^{-n_1}(x), \dots, f^{n-1}(x)$ relativement à $\{U_1, \dots, U_L\}$. Cette application est injective; en effet, les $n$ dernières lettres de $\Phi(w)$ co\"{\i}ncident avec celles de $w$ à, peut-être, une lettre près : celle codée par $f^{n_1}(U_1)$ dans $w$, et qui {n'a pu être transformée qu'en $\infty$ dans $\Phi(w)$ car $f^{n_1}(U_1)$ est disjoint des $U_i$}; or, ce changement peut être déterminé en regardant si la lettre $U_1$ apparaît dans les $n$ premières lettres de $\Phi(w)$. Par suite, on obtient
\[
\sharp \cA_{n}(U_1, f^{n_1}(U_1), U_2, \dots, U_L) \leq \sharp \cA_{n+n_1}(U_1, \dots, U_L)
\]
et ainsi, on a $h^{loc}_{pol}(f;x_1, f^{n_1}(x_{1}), x_2, \dots, x_{L}) \leq h^{loc}_{pol}(f;x_{1}, \dots, x_{L})$ comme promis.\linebreak Remarquons que, par récurrence, cette formule prouve, de façon plus générale, que l'entropie polynomiale d'un ensemble fini $\cF$ ne change pas lorsqu'on ajoute à $\cF$ un nombre fini de points qui sont dans les orbites des points de $\cF$.

Prouvons maintenant le lemme. Soient $x_1, \dots, x_L$ des points de $X \setminus \{\infty\}$. Avec ce qu'on a déjà vu, on peut supposer que les $x_j$ sont deux à deux disjoints. Une récurrence immédiate ramène la preuve du lemme à la formule plus simple
\[
h^{loc}_{pol}(f;x_{1}, \dots, x_{L})=h^{loc}_{pol}(f;f(x_{1}), x_2, \dots, x_{L}),
\]
{qu'il nous reste donc à démontrer}. Choisissons des voisinages deux à deux disjoints $U_1, \dots U_L$ de $x_1, \dots, x_{L}$ tels que $f(U_1)$ soit aussi disjoints des $U_1, \dots, U_L$. Une preuve similaire à celle recourant
à l'application $\Phi$ ci-dessus assure que
\[
\sharp \cA_{n}(f(U_1), U_2, \dots, U_L) \leq \sharp \cA_{n+1}(U_1, \dots, U_L).
\]
Ainsi, $h^{loc}_{pol}(f;f(x_{1}), \dots, x_{L}) \leq h^{loc}_{pol}(f;x_{1}, x_2, \dots, x_{L})$. L'inégalité réciproque est aussi vraie en appliquant $f^{-1}$ (noter bien que toutes les quantités précédemment introduites ne changent pas en substituant $f^{-1}$ à $f$).
\end{proof}

\subsection{Entropie polynomiale et ensemble singulier}\label{ss.singular-set}

On dit que les ensembles $U_{1}, \dots, U_{L}$ de $X \setminus \{\infty\}$ sont \emph{mutuellement singuliers} si pour tout entier $M$, il existe un point $x$ et des instants $n_{1}, \dots, n_{L}$ tels que $f^{n_{i}} (x) \in U_{i}$ pour tout $i$ et, pour tout $j \neq i$, $\left| n_{i}-n_{j} \right|> M$. Un ensemble fini de points $\cS = \{x_{1}, \dots, x_{L} \}$ de $X \setminus \{\infty\}$ est dit \emph{singulier} si, pour tous voisinages $U_{1}, \dots U_{L}$ des $x_1, \dots, x_L$, les $U_i$ sont mutuellement singuliers. Le lecteur pourra vérifier que si $U_{1}, \dots, U_{L}$ sont des compacts de $X$ mutuellement singuliers, alors on peut trouver {un ensemble singulier constitué de} $L$ points $x_{1}, \dots, x_{L}$ avec $x_{i} \in U_{i}$ pour tout $i$. Notons que, {par définition}, un {singleton} est toujours singulier. La proposition suivante montre que l'entropie polynomiale provient toujours des ensembles singuliers.

\begin{prop}\label{prop.singulier}
Soit $\cS$ {une partie finie} de $X \setminus \{\infty\}$. Alors
\[
h^{loc}_{pol}(f; \cS) =\max\{h^{loc}_{pol}(f; \cS'), \cS' \subset \cS \text{ et } \cS' \text{ est singulier}\}.
\]
\end{prop}

Voici deux conséquences assez simples de cette proposition. D'abord, l'entropie polynomiale d'un ensemble fini est plus petit que le cardinal maximal d'un sous-ensemble singulier (cf. la troisième remarque après la proposition~\ref{prop.localization}). Ensuite, dans la formule de localisation
\[
h_{pol}(f) = \sup\{h^{loc}_{pol}(f;x_{1}, \dots, x_{L})\},
\]
de la proposition~\ref{prop.localization}, on peut restreindre le supremum à l'entropie polynomiale des ensembles finis singuliers et, en particulier, $h_{pol}(f)$ est majoré par le cardinal maximal d'un ensemble singulier.


\begin{exam*}
Considérons, comme dans la section~\ref{ss.coding}, {l'application linéaire} $A : (x,y) \mapsto (2x, y/2)$. Notons $\hat A$ la fonction induite par $A$ sur l'espace quotient de $\bbR^2 \cup\{\infty\}$ obtenu en identifiant $0$ et $\infty$; de sorte que $\hat A$ admet un seul point errant. Les ensembles singuliers sont exactement les paires de points constituées d'un point de l'axe des abscisses, différent de l'origine, et d'un point de l'axe des ordonnées, différent de l'origine. Ainsi, on voit que l'entropie polynomiale de $\hat A$ est inférieure à $2$. Réciproquement, dans la section~\ref{ss.coding}, on a trouvé deux disques $Y_1, Y_2$ dont l'entropie polynomiale valait précisément $2$, donc d'après le lemme~\ref{lem.inequality}, l'entropie polynomiale de $\hat A$ est exactement égale à $2$. L'entropie polynomiale de $A$ est aussi égale à $2$ : en effet, on peut facilement étendre les résultats de ce papier au cas où l'ensemble non-errant est fini, la seule preuve où il y a quelque chose de plus à vérifier est celle de la proposition~\ref{prop.localization}. Cette extension pourrait s'appliquer, par exemple, au calcul de l'entropie polynomiale des systèmes dynamiques de Morse--Smale.
\end{exam*}

\begin{proof}[Preuve de la proposition]
Soit $\cS$ un sous-ensemble fini de $X \setminus \{\infty\}$. Supposons que $\cS$ ne soit pas un ensemble singulier, nous allons montrer qu'il existe un sous-ensemble strict $\cS'$ de $\cS$ dont l'entropie polynomiale égale celle de $\cS$. Par récurrence sur la taille de $\cS$ (rappelons qu'un singleton est toujours un ensemble singulier), la proposition sera alors démontrée.

Supposons que $\cS=\{x_1, \dots, x_L\}$ ne soit pas un ensemble singulier. Il existe alors une famille $\cF$ de voisinages errants et deux à deux disjoints $\{U_1, \dots, U_L\}$ des points $x_i$, qui ne sont pas mutuellement singuliers. Soit $n>0$ un entier. On recherche un sous-ensemble strict de $\cF$ de même entropie polynomiale que $\cF$. Comme dans la preuve du lemme~\ref{lem.inequality}, pour tout sous-ensemble $\cF'$ de $\cF$, on note $\cA_n(\cF;\cF')$ l'ensemble des mots de $\cA_n(\cF)$ qui utilisent exactement les lettres de l'alphabet $\cF'\cup\{\infty\}$. Remarquons que puisque les $U_i$ sont errants, toutes les lettres de cet alphabet à l'exception de $\infty$ apparaissent exactement une fois {dans chaque mot $\underline{w} \in \cA_n(\cF;\cF')$}. On obtient la partition :
\[
\cA_n(\cF) = \bigcup_{\cF' \subset \cF} \cA_n(\cF; \cF').
\]
Comme les $U_i$ ne sont pas mutuellement singuliers, il existe un entier $M$ tel que, si une orbite rencontre tous les ensembles $U_i$, alors les instants de passages dans deux d'entre eux diffèrent de moins de $M$. Pour tous $i \neq j \in \{1, \dots, L\}$, soit $\cA_n(\{i,j\})$ l'ensemble des éléments $\underline{w}$ de $\cA_n(\cF;\cF)$ dont les lettres $U_i$ et $U_j$ sont à une distance plus petite que $M$. {On a alors}
\[
\cA_n(\cF;\cF) = \bigcup_{\{i,j\}} \cA_n(\{i,j\}) \ \ \text{ et } \ \
\cA_n(\cF) = \bigcup_{\cF' \subsetneq \cF} \cA_n(\cF; \cF') \cup \bigcup_{\{i,j\}} \cA_n(\{i,j\}).
\]
Soit $C = (2^L-1)+L(L-1)/2$. D'après le principe des tiroirs, sur les $C$ ensembles : $\cA_n(\cF;\cF')$ pour $\cF' \subsetneq \cF$ et $\cA_n(\{i,j\})$ pour $i\neq j$ qui apparaissent dans le membre de droite de la dernière égalité, {l'un au moins} est de cardinal {au moins égal à} à $\frac{\sharp \cA_n(\cF)}{C}$. De plus, l'ensemble des indices est fini et indépendant de $n$. Ainsi, l'entropie polynomiale de $\cF$ provient d'au moins l'un de ces $C$ ensembles; c'est-à-dire qu'au moins un {des cas de figure suivants} a lieu :
\begin{enumerate}
\item \label{item1} il existe un sous-ensemble strict $\cF'$ de $\cF$ tel que :
\begin{equation}\label{eq1}
h_{pol}(\cF)= \limsup_{n \to +\infty} \frac{\log(\sharp\cA_n(\cF;\cF'))}{\log(n)};
\end{equation}

\item \label{item2} il existe une paire $\{i,j\} \text{ avec } i \neq j$ telle que :
\begin{equation}\label{eq2}
h_{pol}(\cF) = \limsup_{n \to +\infty} \frac{\log(\sharp\cA_n(\{i,j\}))}{\log(n)}.
\end{equation}
\end{enumerate}

Examinons la première possibilité. Grâce à l'inclusion suivante (vraie pour tout~$n$) :
\[
\cA_n(\cF;\cF') \subset \cA_n(\cF')
\]
et à l'égalité~\eqref{eq1}, on a $h_{pol}(\cF) \leq h_{pol}(\cF')$. En fait, par décroissance, il y a égalité : $h_{pol}(\cF) = h_{pol}(\cF')$.

Examinons la seconde possibilité. Soit $n>0$ un entier. Si $\underline{w}$ est un mot de $\cA_n(\{i,j\})$, notons $\underline{w}'$ le mot obtenu à partir de $\underline{w}$ en changeant la lettre \og$U_i$\fg{} de $\underline{w}$ (qui apparaît exactement une fois) en \og$\infty$\fg. Comme les $U_k$ sont deux à deux disjoints, $\underline{w}'$ est un mot de $\cA_n(\cF')$ où $\cF'=\cF\setminus\{U_i\}$. De plus, le mot $\underline{w}$ contient les lettres \og$U_j$\fg{} et \og$U_i$\fg{} à une distance plus petite que $M$; par conséquent, $\underline{w}'$ admet au plus ${2M}$ antécédents par l'application $\underline{w} \mapsto \underline{w}'$ qui supprime la lettre \og$U_i$\fg. Cette constatation débouche sur l'inégalité
\[
\cA_n(\{(i,j\}) \leq 2M \sharp \cA_n(\cF').
\]
En faisant tendre $n$ vers $+\infty$ et en recourant à l'égalité~\eqref{eq2}, on conclut comme dans le premier cas que $h_{pol}(\cF) \leq h_{pol}(\cF')$ et, par décroissance, $h_{pol}(\cF) = h_{pol}(\cF')$.

Pour terminer la preuve de la proposition, choisissons pour chaque $i$ une suite décroissante $(V_i(k))_{k \geq 0}$ avec $V_i(0) = U_i$ qui forme une base de voisinages de $x_i$. Considérons maintenant un entier $k$. Comme les $U_i$ ne sont pas mutuellement singuliers, les $V_i(k)$ ne le sont pas non plus. On peut donc appliquer ce que l'on a déjà vu pour les $U_i$ et trouver un sous-ensemble strict $I(k)$ de $\{1, \dots, L\}$ tel que : $ h_{pol}(\{V_i{(k)}, i\in I(k)\}) = h_{pol}(\{V_1(k), \dots, V_L(k)\})$. Comme il n'y a qu'un nombre fini de sous-ensembles de $\{1, \dots, L\}$, l'un d'eux (notons-le $I$) apparaît une infinité de fois dans la suite des $(I_k)_{k\geq 0}$. En couplant cela avec la définition de l'entropie polynomiale, on aboutit à l'égalité
\[
h^{loc}_{pol}(\{x_i, i=1 \dots L\}) = h^{loc}_{pol}(\{x_i, i\in I\}).
\]
Notre proposition est démontrée !
\end{proof}


\section{Entropie polynomiale des homéomorphismes de Brouwer}\label{s.brouwer}

Cette section sera consacrée à montrer le théorème~\ref{theo}. Pour cela, nous commencerons par déterminer l'entropie polynomiale d'une translation (qui se calcule aisément). Ensuite, on montrera que si un homéomorphisme de Brouwer $f$ n'est pas conjugué à une translation, son entropie polynomiale est {au moins égale à} $2$, {ce qui sera} une conséquence d'un résultat classique sur la dynamique des homéomorphismes de Brouwer. Pour finir, on construira pour tout $\alpha \in [2,+\infty]$ un homéomorphisme de Brouwer $f_\alpha$ d'entropie polynomiale $\alpha$. Si $\alpha< +\infty$, l'homéomorphisme $f_{\alpha}$ sera obtenu en \og recollant\fg{} $\lceil \alpha\rceil$ translations. Pour les recollements les plus simples, la construction donnera toujours une entropie polynomiale égale à $2$ quel que soit le nombre de translations; ainsi, les recollements devront être soigneusement \og tordus\fg{} pour disposer d'un nombre assez grand de temps de transition indépendants entre des voisinages d'un ensemble fini de points singuliers. L'homéomorphisme $f_{+\infty}$ peut s'obtenir en \og accolant\fg{} $f_1, f_2, \dots$ ou bien alors, en généralisant la méthode de construction pour le cas fini en recollant une infinité de translations; les détails sont laissés à la récréation du lecteur.

%\bigskip

Lançons-nous dès à présent dans la preuve des deux premiers points du théorème. Soit $f$ un homéomorphisme de Brouwer. Supposons que $f$ soit une translation. Soit $Y$ un compact non vide du plan. Aussi grand que $Y$ soit, il existe un entier $M$ tel que si un point $x$ admet deux itérés $f^{n_1}(x)$ et $f^{n_2}(x)$ qui tombent dans $Y$, alors les temps de passage $n_1$ et $n_2$ diffèrent de moins de $M$ : $|n_1-n_2|\leq M$. Il s'ensuit immédiatement que $h_{pol}(f;Y) = 1$. D'après la proposition~\ref{prop.localization}, on conclut que $h_{pol}(f) = 1\dots$, ce qui est l'implication réciproque du premier point du théorème.

Supposons maintenant que $f$ n'est pas conjugué à un translation. Un résultat classique de théorie des homéomorphismes de Brouwer donne l'existence de deux points $x,y$ tels que $\{x,y\}$ soit une paire singulière au sens où nous l'avons définie dans la section~\ref{ss.singular-set} (cf. par exemple~\cite{Nak95b}). L'argument est probablement dû à Kerekjarto; en voici un résumé : s'il n'existe pas de paire singulière, l'espace des orbites $\bbR^2/f$ est séparé, le {passage au quotient} est un revêtement; ainsi, l'espace quotient est une surface de groupe fondamental $\bbZ$; la classification des surfaces nous apprend que c'est un anneau infini, ce qui nous fournit donc une conjugaison entre $f$ et une translation. Soient $U,V$ deux voisinages compacts errants de $x$ et $y$, soit $A$ l'ensemble des entiers $n$ tels que $f^n(U)$ intersecte $V$. Premièrement, par définition des ensembles singuliers, $A$ contient une infinité d'entiers (positifs ou négatifs). \`A permutation près de $U$ et $V$, on peut supposer que $A$ n'est pas borné supérieurement et contient une infinité d'entiers positifs. De plus, une propriété des homéomorphismes de Brouwer affirme que l'ensemble $A$ est un intervalle de $\mathbb Z$ (c'est un cas particulier du lemme de Franks, se reporter par exemple à~\cite[Lemma~7]{Ler99}). On en déduit l'existence d'un entier $M$ tel que $A$ contienne $\{M,M+1, \dots \}$. On peut alors conclure, comme dans l'exemple de la section~\ref{ss.coding}, que $h_{pol}(f;\{U,V\}) = 2$. {On en déduit $h_{pol}(f) \geq 2$ grâce au lemme~\ref{lem.inequality} ou à la proposition~\ref{prop.localization}}. On a donc prouvé à la fois le point $2$ du théorème et l'implication directe du premier point.

\subsection{Construction d'homéomorphisme de Brouwer par recollement de translations}\label{ss.recollement}

Les dernières pages de ce papier sont consacrées à la construction d'homéomorphismes de Brouwer d'entropie polynomiale plus élevée (point~\eqref{1.1.3} du théorème~\ref{theo}). {Nous rappelons d'abord comment recoller plusieurs translations ensemble pour construire des homéomorphismes de Brouwer plus complexes.} Par souci de simplicité, nous nous restreindrons à la construction d'homéomorphismes qui sont {le temps un d'un flot} (ce qui sera suffisant pour obtenir toutes les valeurs possibles d'entropie polynomiale). Le lecteur avide de plus de détails pourra se reporter à~\cite{Nak95,BegLer03}, et à la figure~6,~162 de la seconde référence pour une intuition géométrique sur cette construction.

Soit $L \geq 2$. Considérons $L$ copies $P_{1}, \dots, P_{L}$ du plan $\bbR^2$ et notons $O_{k}$ le demi-plan supérieur $\{y>0\}$ de $P_{k}$. Pour $k=1, \dots, L-1$, {considérons une application \og changement de cartes\fg{} $\Phi_{k,k+1}: O_{k} \to O_{k+1}$ de la forme suivante} :
\[
(x,y) \longmapsto (x+\varphi_{k,k+1}(y), y)
\]
où $\varphi_{k,k+1}$ est une fonction continue de $\intervalleoo{0}{+\infty}$ dans $\bbR$ qui admet $-\infty$ comme limite en $0^+$. Soit $P$ l'espace quotient :
\[
\cup P_{k} / \sim
\]
où $\sim$ désigne la relation d'équivalence qui identifie {chaque point} $(x,y)$ de $O_k$ avec le point $\Phi_{k,k+1}(x,y)$ de $O_{k+1}$. Nous laissons au lecteur le soin de vérifier que $P$ est une surface séparée, non compacte et simplement connexe; $P$ est donc homéomorphe au plan.

Soit $T : \cup P_{k} \to \cup P_{k}$ l'application de translation $(x,y) \mapsto (x+1,y)$ sur chacun des $P_{k}$. $T$ commute avec tous les $\Phi_{k,k+1}$ et induit donc un homéomorphisme de Brouwer $f: P \to P$.

{Soient $M$ et $M'$ deux points appartenant respectivement à la frontière de $O_{k}$ dans $P_{k}$ et à la frontière de $O_{k'}$ dans $P_{k'}$, pour $k
\neq k'$. On peut alors trouver un point $z$ de $O_{k}$, arbitrairement proche de $M$, qui a un itéré arbitrairement grand qui est proche de $M'$. On en déduit que la paire $\{M,M'\}$ est singulière. Plus généralement,} les ensembles singuliers de $f$ sont les ensembles $\{M_1, \dots, M_L\}$ où $M_k$ est un point de la frontière du demi-plan $O_{k}$ de $P_k$ et leurs sous-ensembles. D'après les propositions~\ref{prop.localization} et~\ref{prop.singulier}, l'entropie polynomiale est le supremum de l'entropie des ensembles singuliers finis. Comme l'entropie polynomiale de $L$ points est {au plus égale} à $L$, on en déduit que $h_{pol}(f_{\alpha}) \leq L$. En particulier, si $L=2$ l'entropie polynomiale est exactement égale à $2$ quel que soit le changement de carte $\Phi_{1,2}$.

\subsection{Construction de \texorpdfstring{$f_\alpha$}{falpha}}

Prenons $L \geq 3$ et $\alpha \in \intervalleof{L-1}{L}$ quelconques. Nous allons maintenant préciser quelles fonctions de changement de cartes $\Phi_{k,k+1}$ nous allons choisir pour construire l'homéomorphisme de Brouwer $f_\alpha$ d'entropie polynomiale $\alpha$ (cf. la figure~\ref{fig.graphs} pour voir à quoi ressemblent les fonctions $\varphi_{k,k+1}$).

\begin{enonce*}{Définition de $\varphi_{k,k+1}$}
Posons $\alpha'=\alpha-L+1$, $\alpha'$ appartient à $\intervalleof{0}{1}$.
\begin{itemize}
\item $\varphi_{1,2}$ est négative, croissante, et, pour chaque entier naturel $k_1$, $\varphi_{1,2}$ prend la valeur $-k_1$ sur un intervalle d'intérieur non vide $I_{k_1}$. Cette famille d'intervalles tend vers $0$ quand $k_1$ tend vers $+\infty$. Pour des questions techniques, on supposera que tous ces intervalles sont inclus dans $\intervalleof{0}{\frac{2}{3}}$.

\item Pour tout entier naturel $k_1$, la restriction de $\varphi_{2,3}$ à $I_{k_1}$ croît de $-2k_1$ à $-k_1$ et, pour tout entier {$k_2$ compris entre $k_1$ et $2k_1$}, $\varphi_{2,3}$ prend la valeur $-k_2$ sur un sous-intervalle $I_{k_1,k_2} \subset I_{k_1}$. On appellera tous ces intervalles les marches d'ordre $k_1$ de $\varphi_{2,3}$. Entre deux marches successives, $\varphi_{2,3}$ est monotone.

\item De même, $\varphi_{3,4}$ est croissante sur chaque marche $I_{k_1,k_2}$ d'ordre $k_1$ de $\varphi_{2,3}$ et {prend chaque valeur entière} $-k_3$ comprise entre $-2k_1$ et $-k_1$ sur un sous-intervalle non vide $I_{k_1,k_2,k_3}$ de $I_{k_1,k_2}$.

\item Et ainsi de suite, jusqu'à $\varphi_{L-1,L}$ : sur chacune des marches d'ordre $k_1$ de $\varphi_{L-2,L-1}$, l'application $\varphi_{L-1,L}$ est croissante et prend toutes les valeurs entières $-k_{L-1}$ entre $-k_1-k_1^{\alpha'}$ et $-k_1$ sur un sous-intervalle $I_{k_1,\dots, k_{L-1}}$ de $I_{k_1,\dots, k_{L-2}}$. Toutes ces fonctions sont monotones entre deux marches successives.
\end{itemize}
\end{enonce*}

\begin{enonce*}[remark]{N.B}
On peut obtenir toutes ces propriétés en prenant des fonctions $\varphi_{k,k+1}$ de classe $C^\infty$; lors, le plan topologique $P$ est muni d'une structure $C^\infty$ pour laquelle $f_\alpha$ est un $C^\infty$-difféomorphisme. Par unicité de la structure $C^\infty$ du plan, on peut transporter $f_\alpha$ en un $C^\infty$-difféomorphisme du plan usuel $\bbR^2$, l'application obtenue est alors le résultat au temps $t=1$ de l'intégration d'un champ de vecteur de classe $C^\infty$ sur le plan.
\end{enonce*}

\clearpage\begin{figure}[p]
\centering\includegraphics[width=245pt]{assets/gluing-staircase.png}\par\medskip
\includegraphics[width=188pt]{assets/gluing-first.png}
\hspace{1cm}
\includegraphics[width=197pt]{assets/gluing-last.png}
\caption{Graphes des fonctions de recollement $\varphi_{\ell,\ell+1}$\label{fig.graphs}}
\end{figure}\clearpage

\subsection{Entropie polynomiale de \texorpdfstring{$f_{\alpha}$}{falpha}}

Le résultat suivant est la clef qui va permettre le calcul de l'entropie polynomiale de $f_\alpha$.

\begin{lemm}\label{lemm.entropy-estimate}
Soit $U_{i}$ l'ensemble $\left[-\frac{2}{3}\,;\frac{2}{3}\right]^2$ dans le plan $P_{i}$. Alors
\[
h_{pol}(U_{1}, \dots, U_{L}) = \alpha.
\]
\end{lemm}

Avant de nous lancer dans la preuve du lemme, voici le déroulé de la suite de la démonstration : pourquoi a-t-on alors bien $h_{pol}(f_\alpha)=\alpha$ ? Par union (point~\eqref{2.2.2} du lemme~\ref{lemm.add-mon}), on a $h_{pol}(U_{1}, \dots, U_{L}) = h_{pol}(U_{1} \cup \dots \cup U_{L})$. Puis, en vertu du principe de localisation (proposition~\ref{prop.localization}), on obtient le minorant voulu :
\[
h_{pol}(f_{\alpha}) \geq \alpha.
\]
Pour obtenir le majorant, on recourt à la formule qui relie l'entropie polynomiale à l'entropie locale des ensembles singuliers de $f_{\alpha}$, que l'on a décrits à la fin de la section~\ref{ss.recollement}; on obtient :
\[
h_{pol}(f_{\alpha}) = \sup \left\{h^{loc}_{pol}(f_{\alpha}; x_{1}, \dots, x_{L}), x_{k} \in \partial O_{k} \text{ pour } k=1, \dots, L\right\}.
\]
De plus, le lemme~\ref{lemm.iterates} assure que remplacer dans le calcul de l'entropie un point par un de ses itérés n'a aucune influence sur le résultat. Comme tout point de $\partial O_{k}$ comporte un itéré qui tombe dans {l'intérieur de} $U_{k}$, on peut, en toute généralité, supposer que $x_{k}$ appartient à l'intersection de $\partial O_{k}$ avec l'intérieur de $U_{k}$. Or, par définition, l'entropie locale aux points $x_{1}, \dots, x_{L}$ est inférieure ou égale à l'entropie de $\{U_{1}, \dots, U_{L}\}$; il en va donc de même pour $h_{pol}(f_{\alpha})$. Cela montre que $\alpha$ est aussi un majorant de l'entropie polynomiale. (Cet argument permet aussi de voir que tout compact dont l'intérieur intersecte chacun des $\partial O_k$ a comme entropie polynomiale $h_{pol}(f_\alpha)$.)

%\bigskip

Il ne nous reste donc plus qu'à montrer le lemme~\ref{lemm.entropy-estimate}.

\begin{proofc}[Preuve du lemme]

\begin{step}[Minoration : $h_{pol}(f_\alpha) \geq \alpha$]
Soit $n>0$ un entier, on veut un minorant du nombre de mots dans $\cA_n (f; \{U_1, \dots, U_L\})$.

Soit $k_1$ un entier naturel plus petit que $\frac{n}{2L}$. Soient $k_2, \dots, k_{L-2}$ des entiers compris entre $k_1+1$ et $2k_1$. Soit $k_{L-1}$ un entier compris entre $k_1+1$ et $k_1+k_1^{\alpha'}$. Finalement, choisissons $y$ dans l'intervalle $I_{k_1,\dots, k_{L-1}}$. En notant $z$ le point du plan $P$ dont les coordonnées dans $P_1$ sont $(0,y)$, le point $z$ est dans $U_1$. Comme $\varphi_{1,2}(y)=-k_1$, $z$ a pour coordonnées $(-k_1, y)$ dans le plan $P_2$, ainsi un itéré $f^{k}(z)$ appartient à $U_2$ si et seulement si $k=k_1$. De même, $\varphi_{2,3}(y)=-k_2$, les coordonnées de $f^{k_1}(z)$ dans $P_3$ sont $(-k_2,0)$; un itéré $f^{k}(f^{k_1}(z))$ est dans $U_3$ si et seulement si $k=k_2$, etc. Soit $k_0$ un entier compris entre $1$ et $k_1$. Comme $k_1$ est plus petit que $\frac{n}{2L}$, l'encodage de l'orbite des $n$ premiers itérés de $f^{-k_0}z$ commence par :
\[
\underbrace{\infty, \dots,\infty}_{k_0 \text{ lettres}}, U_1,
\underbrace{\infty, \dots,\infty}_{(k_1 -1) \text{ lettres}}, U_2,
\underbrace{\infty, \dots,\infty}_{(k_2 -1) \text{ lettres}}, U_3, \dots \dots, U_{L-1},
\underbrace{\infty, \dots, \infty}_{(k_{L-1} -1) \text{ lettres}},U_L.
\]
On voit que deux suites $(k_\ell)_\ell$ distinctes fournissent deux encodages différents. De plus, pour une valeur fixée de $k_1$, on a ${k_1}^{L-2}\lfloor{k_1}^{\alpha'}\rfloor \geq (k_1-1)^{\alpha-1}$ possibilités pour le $(L-1)$-uplet $(k_0, k_2, k_3, \dots, k_{L-1})$. On obtient de la sorte le minorant suivant :
\[
\sharp \cA_n (f; \{U_1, \dots, U_L\}) \geq
\sum_{1 \leq k \leq \frac{n}{2L}} (k-1)^{\alpha-1}.
\]
En comparant par exemple cette somme avec une intégrale, on obtient l'équivalent (quand $n$ tend vers $+\infty$) :
\[
\frac{1}{(2L)^\alpha \alpha} n^\alpha
\]
(On aurait d'ailleurs pu se limiter aux valeurs de $n$ qui sont multiples de $2L$ puisqu'on ne calcule qu'une limite supérieure). Ainsi, $h_{pol}(f; \{U_1, \dots, U_L\}) \geq \alpha$.
\end{step}

\begin{step}[Majoration : $h_{pol}(f_\alpha) \leq \alpha$]
On considère encore un entier $n$, il nous faut un majorant au nombre d'éléments de $\cA_n (f; \{U_1, \dots, U_L\})$. Comme chaque fonction $\varphi_{k,k+1}$ est négative, si l'on excepte l'élément ne comportant que la lettre $\infty$, tout autre mot est de la forme suivante :
\[
{\infty, \dots,\infty}, U_i, {\infty, \dots,\infty}, U_{i+1},
\dots \dots, U_{j}, {\infty, \dots, \infty}
\]
(avec $i \leq j$ dans $\{1, \dots, L\}$) ou d'une forme similaire avec des lettres $U_k$ possiblement dédoublées (un point et son itéré peuvent appartenir tous deux à $U_k$). Si $i>1$ ou $j <L$, {ce mot contient au plus} $L$ sous-mots {maximaux} composés uniquement de la lettre $\infty$, tous de taille bornée par $n$ et la taille du dernier tel sous-mot est déterminé par la taille des autres. Même en prenant en compte les éventuels dédoublements, cela ne fait toujours que (à $i$ et $j$ fixés)
\[
2^{L-1}n^{L-1}
\]
possibilités au maximum, qui est asymptotiquement borné par $n^\alpha$. Il ne reste qu'à estimer le nombre de possibilités quand $i=1$ et $j=L$. Soit $z$ un point et $(x,y)$ ses coordonnées dans $P_1$. Si le point $z$ appartient à $U_1$ et que $f^k(z)$ appartient à $U_2$, alors
\[
-\varphi_{1,2}(y) \in \left[k-\frac{4}{3}\,;k+\frac{4}{3}\right]
\]
En effet, si $x_1$ et $x_2$ sont les premières coordonnées de $z$ pour, respectivement, $P_1$ et $P_2$, alors $x_1, x_2+k \in \intervalle{-\frac{2}{3}}{\frac{2}{3}}$ et $x_2 = x_{1} +\varphi_{1,2}(y)$. La définition de $\varphi_{1,2}$ entraîne alors que $y$ se trouve entre le maximum de $I_{k-2}$ et le minimum de $I_{k+2}$. Par conséquent $-\varphi_{L-1,L}(y)$ tombe dans l'intervalle
\[
[k-2, k+2+(k+2)^{\alpha'}].
\]
Regardons à quoi ressemble un élément de $\cA_n (f; \{U_1, \dots, U_L\})$ :
\[
w={\infty, \dots,\infty}, U_1,
\underbrace{\infty, \dots,\infty}_{k_1-1 \text{ lettres}}, U_{2},
\dots \dots, U_{L-1}
\underbrace{\infty, \dots,\infty}_{k_{L-1}-1 \text{ lettres}}, U_{L}, {\infty, \dots, \infty}.
\]
Il y a au plus $L+1$ sous-mots {maximaux} composés seulement de la lettre $\infty$. Chacun de ces sous-mots est de taille moins que $n$. De plus, l'une des tailles est déterminée par les autres, {et d'après l'estimation sur $-\varphi_{L-1,L}(y)$,} pour $k_1$ fixé, il y a au plus $(k_1+2)^{\alpha'}+5$ possibilités pour $k_{L-1}$. Pour finir, toujours à $k_1$ fixé, {parmi les $L+1$ sous-mots qui ne contiennent que la lettre $\infty$, le second a pour longueur $k_{1}$, l'avant-dernière longueur n'a que $(k_1+2)^{\alpha'}+5$ possibilités, il reste $L-2$ longueurs bornées par $n$ à choisir (la dernière est déterminé par les autres).} Au total, {en tenant compte de toutes les possibilités pour $k_{1}$}, le nombre de possibilités pour de tels mots est borné par :
\[
n^{L-2} \sum_{k=1}^n ((k+2)^{\alpha'}+5) \leq n^{L-2} \left((n+2)^{\alpha'+1} + 5n\right) \sim n^{\alpha}.
\]
On aboutit à $h_{pol}(f; \{U_1, \dots, U_L\}) \leq \alpha$ et tout est accompli !\qed
\end{step}
\let\qed\relax
\end{proofc}

\begin{thebibliography}{MMMMMM}
\bibitem[BL16]{BerLab16} Bernard, Patrick; Labrousse, Clémence An entropic characterization of the flat metrics on the two torus, Geom. Dedicata, Volume 180 (2016), pp. 187-201 \href{https://ahl.centre-mersenne.org/item/AHL_2019__2__39_0/#BerLab16}{Publisher reference}.
\bibitem[BLR03]{BegLer03} Béguin, François; Le Roux, Frédéric Ensemble oscillant d’un homéomorphisme de Brouwer, homéomorphismes de Reeb, Bull. Soc. Math. Fr., Volume 131 (2003) no. 2, pp. 149-210 \href{https://ahl.centre-mersenne.org/item/AHL_2019__2__39_0/#BegLer03}{Publisher reference}.
\bibitem[CK97]{CasKar97} Cassaigne, Julien; Karhumäki, Juhani Toeplitz words, generalized periodicity and periodically iterated morphisms, Eur. J. Comb., Volume 18 (1997) no. 5, pp. 497-510 \href{https://ahl.centre-mersenne.org/item/AHL_2019__2__39_0/#CasKar97}{Publisher reference}.
\bibitem[HLR17]{englishversion} Hauseux, Louis; Le Roux, Frédéric Polynomial entropy of Brouwer homeomorphisms (2017) (\url{https://arxiv.org/abs/1712.01502}) \href{https://ahl.centre-mersenne.org/item/AHL_2019__2__39_0/#englishversion}{Publisher reference}.
\bibitem[Kan18]{Kan18} Kanigowski, Adam Slow entropy for some smooth flows on surfaces, Isr. J. Math., Volume 226 (2018) no. 2, pp. 535-577 \href{https://ahl.centre-mersenne.org/item/AHL_2019__2__39_0/#Kan18}{Publisher reference}.
\bibitem[KT97]{KatTho97} Katok, Anatole; Thouvenot, Jean-Paul Slow entropy type invariants and smooth realization of commuting measure-preserving transformations, Ann. Inst. Henri Poincaré, Probab. Stat., Volume 33 (1997) no. 3, pp. 323-338 \href{https://ahl.centre-mersenne.org/item/AHL_2019__2__39_0/#KatTho97}{Publisher reference}.
\bibitem[Kus67]{Kus67} Kushnirenko, A. G. Metric invariants of entropy type, Usp. Mat. Nauk, Volume 22 (1967) no. 5 (137), pp. 57-65 \href{https://ahl.centre-mersenne.org/item/AHL_2019__2__39_0/#Kus67}{Publisher reference}.
\bibitem[KVW18]{KVW18} Kanigowski, Adam; Vinhage, Kurt; Wei, Daren Kakutani Equivalence of Unipotent Flows (2018) (\url{https://arxiv.org/abs/1805.01501}) \href{https://ahl.centre-mersenne.org/item/AHL_2019__2__39_0/#KVW18}{Publisher reference}.
\bibitem[Lab13]{Lab13} Labrousse, Clémence Polynomial entropy for the circle homeomorphisms and for C 1  nonvanishing vector fields on 𝕋 2  (2013) (\url{https://arxiv.org/abs/1311.0213}) \href{https://ahl.centre-mersenne.org/item/AHL_2019__2__39_0/#Lab13}{Publisher reference}.
\bibitem[LR99]{Ler99} Le Roux, Frédéric Bounded recurrent sets for planar homeomorphisms, Ergodic Theory Dyn. Syst., Volume 19 (1999) no. 4, pp. 1085-1091 \href{https://ahl.centre-mersenne.org/item/AHL_2019__2__39_0/#Ler99}{Publisher reference}.
\bibitem[Mar13]{Mar13} Marco, Jean-Pierre Polynomial entropies and integrable Hamiltonian systems, Regul. Chaotic Dyn., Volume 18 (2013) no. 6, pp. 623-655 \href{https://ahl.centre-mersenne.org/item/AHL_2019__2__39_0/#Mar13}{Publisher reference}.
\bibitem[Nak95a]{Nak95} Nakayama, Hiromichi A non-flowable plane homeomorphism whose non-Hausdorff set consists of two disjoint lines, Houston J. Math., Volume 21 (1995) no. 3, pp. 569-572 \href{https://ahl.centre-mersenne.org/item/AHL_2019__2__39_0/#Nak95}{Publisher reference}.
\bibitem[Nak95b]{Nak95b} Nakayama, Hiromichi On dimensions of non-Hausdorff sets for plane homeomorphisms, J. Math. Soc. Japan, Volume 47 (1995) no. 4, pp. 789-793 \href{https://ahl.centre-mersenne.org/item/AHL_2019__2__39_0/#Nak95b}{Publisher reference}.
\bibitem[Wal82]{Wal82} Walters, Peter An introduction to ergodic theory, Graduate Texts in Mathematics, 79, Springer, 1982, ix+250 pages \href{https://ahl.centre-mersenne.org/item/AHL_2019__2__39_0/#Wal82}{Publisher reference}.
\end{thebibliography}
\end{document}
